\documentclass[11pt]{article}
\usepackage[margin=0.85in]{geometry}
\usepackage{amsmath,amssymb,amsthm}
\usepackage{hyperref}
\newtheorem{theorem}{Theorem}
\newtheorem{lemma}{Lemma}
\title{A finite-certificate proof of the cube-sum permutation conjecture (OEIS A006063)}
\author{Research note for independent audit}
\date{18 October 2026}
\begin{document}
\maketitle

\begin{abstract}
A permutation $p$ of $\{1,\ldots,n\}$ is called a cube-sum
permutation if $i+p_i$ is a perfect cube for every $i$.
OEIS A006063 conjectures that such a permutation exists for every
$n\ge 173$. We reduce the infinite assertion to only $192$ finite
cases ($173\le n\le364$) by an elementary interval-reflection
extension. We supply the complete $192$ integer permutation witnesses
($51{,}552$ entries) and an exact standalone checker. The resulting
argument is a computer-assisted proof candidate, pending independent
checking of the certificate and literature priority.
\end{abstract}

\begin{lemma}[Reflection extension]\label{lem:extension}
Suppose $0\le m<n$ and a cube-sum permutation exists for $m$.
If $m+n+1=c$ is a perfect cube, then one exists for $n$.
\end{lemma}
\begin{proof}
Retain the original permutation on $\{1,\ldots,m\}$ and, for
$m+1\le i\le n$, define $p_i=c-i$. Because $c=m+n+1$,
the map $i\mapsto c-i$ is a bijection of the interval
$\{m+1,\ldots,n\}$ to itself and $i+p_i=c$ for every new index.
Thus the union is a permutation of $\{1,\ldots,n\}$ satisfying
the required condition.
\end{proof}

\begin{lemma}[Every large interval contains a usable cube]\label{lem:cube}
For every integer $n\ge365$, there is a perfect cube $c$
satisfying $n+174\le c\le2n$.
\end{lemma}
\begin{proof}
For $365\le n\le555$, take $c=9^3=729$.
Now suppose $n\ge556$ and let $k$ be the least integer with
$k^3\ge n+174$. Then $k\ge10$ and
$(k-1)^3<n+174$, hence, by integrality,
$n\ge(k-1)^3-173$. For every integer $k\ge10$,
\[
2((k-1)^3-173)-k^3
=k^3-6k^2+6k-348\ge0.
\]
The last expression is $112$ at $k=10$, and its successive
difference is $3k^2-9k+1>0$ for $k\ge10$.
Consequently $k^3\le2n$. Taking $c=k^3$ proves the claim.
\end{proof}

\begin{lemma}[Finite base certificate]\label{lem:finite}
For every integer $173\le n\le364$, a cube-sum permutation exists.
\end{lemma}
\begin{proof}
The accompanying file
\texttt{certificates/base\_permutations.json} is an ordered array
of $192$ explicit permutations, one for each $n=173,\ldots,364$.
Each row $p$ is independently certified by the elementary
integer-only checks
\[
|p|=n,\qquad
\operatorname{sort}(p)=(1,\ldots,n),\qquad
\forall\, i\in\{1,\ldots,n\}\ \exists\,k\in\mathbb Z_{\ge2}:
\ i+p_i=k^3.
\]
The file contains $51{,}552$ integer entries in total.
These checks are implemented without floating-point arithmetic
or external packages in \texttt{verify.py}.
The separate script \texttt{independent\_match.py}
regenerates independent witnesses using augmenting paths,
without reading the archived certificate.
\end{proof}

\begin{theorem}\label{thm:main}
For every integer $n\ge173$ there exists a permutation
$p_1,\ldots,p_n$ of $1,\ldots,n$ such that
$i+p_i$ is a perfect cube for every $1\le i\le n$.
Equivalently, the conjectured tail of OEIS A006063 is complete.
\end{theorem}
\begin{proof}
Use strong induction on $n$. The base cases
$173\le n\le364$ follow from Lemma~\ref{lem:finite}.
For $n\ge365$, Lemma~\ref{lem:cube} provides
$n+174\le c\le2n$ with $c$ a perfect cube.
Set $m=c-n-1$. Then $173\le m<n$.
By the induction hypothesis, $m$ has a cube-sum permutation.
Since $m+n+1=c$, Lemma~\ref{lem:extension} extends it to $n$.
\end{proof}

\paragraph{Reproduction and scope.}
Run \texttt{python verify.py} and
\texttt{python independent\_match.py} in the
\texttt{papers/A006063} directory. Both use only the Python
standard library. The mathematical induction proves an unbounded
result; numerical checks above $364$ are optional implementation
smoke tests, not premises of the proof. The present note does not
claim literature priority, peer review, or knowledge of the
number of cube-sum permutations for a given size.

\paragraph{Source.}
N.~J.~A.~Sloane et al.,
\href{https://oeis.org/A006063}{OEIS A006063},
including Robert Israel's 2018 existence conjecture and
finite verification to $1000$.
\end{document}
